% =====================================================================
%  共享版式 —— 所有合集共用（由 tools/build.py 拷进构建目录）
%  改这里会影响全部合集；改完重新 build 即可
% =====================================================================

\usepackage{amsmath,amssymb}
\usepackage{geometry}
\geometry{a4paper,top=2.5cm,bottom=2.3cm,left=2.3cm,right=2.3cm,headsep=0.55cm}
\usepackage{enumitem}
\usepackage{needspace}
\usepackage{fancyhdr}
\usepackage{array}
\usepackage{booktabs}
\usepackage[hidelinks]{hyperref}

\linespread{1.28}
\setlength{\parindent}{2em}
\setlength{\abovedisplayskip}{5pt plus 2pt minus 2pt}
\setlength{\belowdisplayskip}{5pt plus 2pt minus 2pt}
\setlength{\abovedisplayshortskip}{3pt plus 1pt}
\setlength{\belowdisplayshortskip}{3pt plus 1pt}

% 书级变量由构建脚本写入 main.tex
\providecommand{\schoolname}{浙江大学}
\providecommand{\coursename}{《线性代数（甲）》}
\providecommand{\booktitle}{历年试题合集}
\newcommand{\currentpaper}{}

% ---------- 题目环境 ----------
% \begin{problem}{15 分}{9cm} 题干 \end{problem}
% 第一参数：分值（留空则只显示题号）；第二参数：题后解答留白高度
\newcounter{problem}
\newcommand{\resetproblem}{\setcounter{problem}{0}}
\newcommand{\problemspace}{0pt}
\newenvironment{problem}[2]{%
  \par\addvspace{1.0em}%
  \needspace{\dimexpr#2+2.4\baselineskip\relax}%
  \renewcommand{\problemspace}{#2}%
  \stepcounter{problem}%
  \noindent\textbf{\theproblem.}%
  \ifx\relax#1\relax\else\hspace{0.25em}\textbf{（#1）}\fi
  \hspace{0.4em}\ignorespaces
}{%
  \par\nobreak\vspace{\problemspace}\par
}

% 填空作答横线：\blank 或 \blank[5cm]
\newcommand{\blank}[1][2.4cm]{\underline{\hspace{#1}}}

% 卷内分块标题，如 \blockhead{填空题（85 分）}
\newcommand{\blockhead}[1]{%
  \par\addvspace{1.3em}\nopagebreak[3]%
  \noindent{\large\bfseries #1}%
  \par\nopagebreak\addvspace{0.45em}%
}

% 每套试卷的卷头：\paperhead{学年}{学期}{考试}{目录条目}
\newcommand{\paperhead}[4]{%
  \clearpage
  \phantomsection
  \addcontentsline{toc}{section}{#4}%
  \renewcommand{\currentpaper}{#1\ #2}%
  \markboth{#4}{#4}%
  \begin{center}
    {\Large\bfseries\ziju{0.12} \schoolname}\\[0.45em]
    {\large \coursename}\\[0.6em]
    {\Large\bfseries #1 学年\quad #2\quad #3考试试卷}\\[0.65em]
    \rule{\linewidth}{1.1pt}
  \end{center}
  \vspace{0.4em}
  \resetproblem
}

\pagestyle{fancy}
\fancyhf{}
\fancyhead[L]{\small \schoolname\booktitle}
\fancyhead[R]{\small\currentpaper}
\fancyfoot[C]{\small\thepage}
\renewcommand{\headrulewidth}{0.4pt}
\fancypagestyle{plain}{\fancyhf{}\fancyfoot[C]{\small\thepage}\renewcommand{\headrulewidth}{0pt}}
